\label{sec:appendix}

\subsection*{Detailed Genre Comparison}

\begin{table}[h]
\centering
\caption{Genre comparison across three analytical dimensions (detailed)}
\label{tab:genre_comparison_appendix}
\small
\resizebox{\columnwidth}{!}{%
\begin{tabular}{lcccccc}
\toprule
\textbf{Layer} & \textbf{Canon.} & \textbf{Pulp} & \textbf{LLM} & \textbf{Can-Pulp} & \textbf{Can-LLM} & \textbf{$p$} \\
\midrule
Cultural (L4)     & 3.96 & 3.83 & 2.55 & +0.13 & \textbf{+1.41} & $<$0.001*** \\
                  & (0.41) & (0.36) & (0.62) & & & \\
\midrule
Emotional (L5)    & 4.15 & 4.04 & 3.36 & +0.11 & \textbf{+0.79} & $<$0.001*** \\
                  & (0.31) & (0.24) & (0.59) & & & \\
\midrule
Existential (L6)  & 3.95 & 3.68 & 2.59 & +0.27 & \textbf{+1.36} & $<$0.001*** \\
                  & (0.57) & (0.56) & (0.57) & & & \\
\midrule
\textbf{Average} & \textbf{4.02} & \textbf{3.85} & \textbf{2.83} & \textbf{+0.17} & \textbf{+1.19} & $<$\textbf{0.001***} \\
\bottomrule
\end{tabular}
}
\vspace{0.1cm}

\small
\textit{Notes:} Mean scores with standard deviations in parentheses. Can-Pulp/Can-LLM = score differences. Statistical significance: ***$p<0.001$. Effect sizes (Cohen's $d$) for Canonical vs LLM: L4 $d$=2.68, L5 $d$=1.68, L6 $d$=2.40 (all very large effects).
\end{table}

\subsection*{Iterative Evaluation Protocol}

Table~\ref{tab:eval_protocol} summarizes the five-round iterative evaluation protocol. All three layers share the same round structure, with layer-specific adaptations noted.

\begin{table}[h]
\centering
\caption{Five-round iterative evaluation protocol}
\label{tab:eval_protocol}
\small
\resizebox{\columnwidth}{!}{%
\begin{tabular}{clp{5.5cm}}
\toprule
\textbf{Round} & \textbf{Focus} & \textbf{Task} \\
\midrule
R1 & Initial Assessment & Extract layer-relevant content; score all 4 dimensions (1.0--5.0) with confidence ratings and textual evidence \\
\midrule
R2 & Bias Check & \textbf{L4}: Hallucination check (fabricated cultural/historical details). \textbf{L5/L6}: Projection bias check (imposing Western emotional norms or existentialist frameworks) \\
\midrule
R3 & Layer Boundary & Verify no cross-layer contamination (e.g., L5 must not evaluate cultural power structures [L4] or philosophical claims [L6]) \\
\midrule
R4 & Evidence Sufficiency & Assess evidence quality; adjust \textit{confidence} (not scores) if evidence is ambiguous or sparse \\
\midrule
R5 & Final Consolidation & Confirm or adjust final scores; convergence criterion: $|\Delta\text{score}| < 0.3$ from R4 to R5 \\
\bottomrule
\end{tabular}
}
\end{table}

The independent validator receives the iterative evaluator's five-round output and performs single-pass cross-verification covering: (1)~independent scoring, (2)~projection bias detection, (3)~hallucination detection, (4)~reasoning quality assessment, and (5)~confidence calibration. Agreement between evaluator and validator is measured per dimension, with discrepancies $>$0.5 points flagged for review.

\subsection*{Dimension Scoring Criteria}

Each dimension is scored on a 1.0--5.0 scale. Table~\ref{tab:scoring_rubric} provides the scoring rubric applied uniformly across all layers.

\begin{table}[h]
\centering
\caption{Scoring rubric for dimension evaluation}
\label{tab:scoring_rubric}
\small
\begin{tabular}{cp{6.5cm}}
\toprule
\textbf{Score} & \textbf{Criteria} \\
\midrule
4.5--5.0 & \textbf{Exceptional}: Rich, multi-layered representation with extensive textual evidence; sophisticated treatment demonstrating mastery \\
4.0--4.5 & \textbf{Strong}: Substantial representation with clear evidence; demonstrates depth and nuance \\
3.0--3.5 & \textbf{Moderate}: Some representation present but limited in scope or depth; evidence adequate but not extensive \\
2.0--2.5 & \textbf{Weak}: Minimal representation; surface-level treatment with sparse evidence \\
1.0--1.5 & \textbf{Absent/Negligible}: Dimension effectively absent or trivially treated; no meaningful evidence \\
\bottomrule
\end{tabular}
\end{table}

\subsection*{Prompt Template (Representative Example)}

The following presents a condensed version of the content-limit system prompt for Layer~5 (Emotional-Psychological), illustrating the prompt structure shared across all three layers. Layer~4 and Layer~6 follow the same template with layer-specific dimension definitions and theoretical frameworks.

\begin{quote}
\small
\ttfamily

\textbf{System Prompt (condensed):}

You are an expert literary-emotional analyst evaluating how texts represent emotions, psychological states, and affective experiences.

\textbf{Core Principles:}\\
1. Evidence Sovereignty: All scores must cite specific textual evidence.\\
2. Strict Layer Boundary: Evaluate ONLY emotional-psychological content. Do NOT evaluate cultural power structures (Layer~4) or philosophical themes (Layer~6).\\
3. Projection Awareness: Avoid imposing assumptions about ``how emotions should be expressed.'' Restrained expression $\neq$ shallow.\\
4. Dimension Independence: Evaluate each dimension independently.

\textbf{Mode: content-limit}\\
Evaluate based SOLELY on textual evidence. Do NOT use prior knowledge about literary works or authors.

\textbf{Dimensions:}\\
1. AC (Affective Complexity) -- Multiplicity, contradiction, and evolution of emotional states. [Sedgwick, Berlant, Ngai]\\
2. PI (Psychological Interiority) -- Depth and accessibility of characters' inner psychological worlds. [Cohn, Bruner, Bakhtin]\\
3. EG (Emotional Granularity) -- Precision and differentiation of emotional vocabulary. [Barrett]\\
4. ENC (Emotional-Narrative Coherence) -- Whether emotions arise organically from narrative situation. [New Criticism, James Wood]

\medskip

\textbf{Round 1 User Prompt (condensed):}

ROUND 1/5: EMOTIONAL CONTENT EXTRACTION

Step 1: Extract emotional content (explicit emotion words, psychological states, interior techniques, affective moments).\\
Step 2: Score each dimension (1.0--5.0) with confidence (1--5), reasoning, and evidence strength.

Output: JSON with extracted\_content, scores (per dimension: score, confidence, reasoning, evidence\_strength), overall\_score.

\medskip

\textbf{Round 2 User Prompt (condensed):}

ROUND 2/5: PROJECTION BIAS CHECK

Review Round~1. Ask critically:\\
-- Did I penalize restrained emotional expression as ``shallow''?\\
-- Did I apply Western emotional display norms?\\
-- Did I reward explicit language without checking substance?\\
Revise scores if bias detected; confirm if none found.
\end{quote}

Rounds 3--5 follow the protocol described in Table~\ref{tab:eval_protocol}. The complete prompt templates, including all JSON output schemas and layer-specific dimension definitions for Layers 4 and 6, are available in the project repository.
